% \section{Introdução}
\label{sec_introducao_formulacao_problema}

Neste capítulo é feita a formulação matemática do problema de encontro e acoplamento ou atracação entre duas espaçonaves.
Este problema é conhecido na literatura como \gls{rvdb}.
O \gls{rvdb} é consistente com a proposta dessa pesquisa de mestrado.
A formulação matemática envolve a dinâmica e controle de órbita e atitude das espaçonaves com movimentos orbital no referencial absoluto (\gls{eci}) e relativo, no referencial orbital, \gls{lvlh}.
As espaçonaves são definidas como perseguidora e alvo.
Para o alvo assume-se órbita e atitude conhecidas e nave não ativa.
A nave perseguidora deve encontrá-la e se acoplar com ela de forma controlada, de modo a prevenir colisões e garantir o sucesso da missão.
Por sua vez, a espaçonave perseguidora é uma \gls{etmr}, responsável pelo acoplamento com o alvo.
O movimento relativo refere-se ao movimento entre as espaçonaves, além disso, o sistema de coordenadas orbital \gls{lvlh} é definido na origem do centro de massa do alvo e segue com a órbita.

A formulação do problema de controle aplica-se a ambas representações do movimento no \gls{eci} e \gls{lvlh}, para órbita e atitude, nas versões não-linear e linerarizada das equações do movimento relativo.
Nesta pesquisa assume-se que o lançador injeta a espaçonave perseguidora em uma órbita coplanar à órbita do alvo.
A órbita final que antecede a aproximação de longa distância, deve ser estável e adequada para se ter início às manobras de fase (de longa, curta e de proximidade (final)), todas definidas como navegação relativa.

Terminando as manobras de longa distância, e continuando a abordagem do movimento relativo, passam-se às manobras de curta distância.
Como se assumiu que a espaçonave alvo não é colaborativa, no sentido de não contribuir para o controle ou para o êxito das operações de \gls{rvdb}, para que ocorra o sucesso da missão, é necessário que a atitude das espaçonaves esteja sincronizada.
